\centerline{\bf \Huge Начало комбинаторики I}

\begin{flushright}
        \scriptsize{Никогда не сдаваться... Встать, когда все рухнуло — вот настоящая сила. Наруто}
\end{flushright}

\problem{\textbf{1}} Сколько двузначных чисел можно составить из цифр 1, 2, 3?

\problem{\textbf{2}} Сколькими способами можно выложить в ряд два красных и три синих шарика? Шарики
не отличаются ничем, кроме цвета.

\problem{\textbf{3}} Сколько трёхзначных чисел можно составить из цифр 0, 1, 2, 3, 4 если цифры в записи числа
не могут повторяться?

\problem{\textbf{4}} Анаграмма — это слово (не обязательно осмысленное), полученное из данного слова перестановкой букв. Например, бьорд является анаграммой слова дробь. Сколько анаграмм имеют
слова маг, дед, краб, ирис?

\problem{\textbf{5}} Сигмабой, Тунтунтунсагур и Камерамэн пишут контрольную по математике. Сигмабой может получить четвёрку, тройку или двойку, Тунтунтунсагур — пятёрку, четвёрку или тройку, а Камерамэн — пятёрку или четвёрку. Сколькими способами может завершиться для них контрольная?

\problem{\textbf{6}} На клетчатой бумаге нарисовали квадрат 6 $\times$ 6. Сколько всего квадратов можно выделить?

\problem{\textbf{7}} На клетчатой бумаге нарисовали квадрат 7 $\times$ 7. Сколько всего прямоугольников можно выделить?